\section*{Capítulo \ref{ch:b3:electrode-kinetics} --- Cinética de eletrodo e eletroanálise}

\begin{solution}{exo:b3:electrode-kinetics:1}
A inclinação catódica $2.303RT/\alpha F = \qty{118}{mV}$ dá $\alpha = 59.2/118 = 0.50$; $j_0 = \qty{2.0e-6}{A.cm^{-2}}$,
o intercepto em $\eta = 0$.
\end{solution}

\begin{solution}{exo:b3:electrode-kinetics:2}
$i_0 = \qty{2.0e-4}{A}$; $R_{\mathrm{ct}} = RT/(Fi_0) = 0.02569/\num{2.0e-4} = \qty{128}{\ohm}$.
\end{solution}

\begin{solution}{exo:b3:electrode-kinetics:3}
Para \ce{MgSO4}, $I = \frac12(4c + 4c) = 4c$. A \qty{1.0}{mmol/L}, $I = \qty{4.0}{mol.m^{-3}}$ e $\kappa^{-1} =
\qty{0.96}{nm} \times \sqrt{100/4} = \qty{4.8}{nm}$; a \qty{0.10}{mol/L}, $I = \qty{400}{mol.m^{-3}}$, $\kappa^{-1} = \qty{0.48}{nm}$.
\end{solution}

\begin{solution}{exo:b3:electrode-kinetics:4}
$i_{\mathrm c} = C_{\mathrm{dl}}Av = \num{20e-6} \times 0.0707 \times 0.1 = \qty{0.14}{\micro A}$, menos de 1\,\% do pico. A
\qty{10}{V/s}, ela vale \qty{14}{\micro A}, enquanto o pico só cresce até $19\sqrt{100} = \qty{190}{\micro A}$:
a corrente capacitiva cresce como $v$, a faradaica como $\sqrt v$.
\end{solution}

\begin{solution}{exo:b3:electrode-kinetics:5}
$D = \pi\bigl(i\sqrt t/nFAc^*\bigr)^2 = \pi(\num{12.2e-6}/(96\,485 \times 0.0707 \times \num{1.00e-6}))^2 =
\qty{1.0e-5}{cm^2.s^{-1}}$.
\end{solution}

\begin{solution}{exo:b3:electrode-kinetics:6}
$FAc^* = \qty{1.364e-2}{C.cm^{-1}}$, $F/RT = \qty{38.92}{V^{-1}}$. A \qty{50}{mV/s}: $\sqrt{38.92 \times 0.05 \times \num{7.0e-6}} =
\num{3.69e-3}$, $i_{\mathrm p} = 0.4463 \times \num{1.364e-2} \times \num{3.69e-3} = \qty{22.5}{\micro A}$; a \qty{500}{mV/s}, $\sqrt{10}$ vezes
mais, \qty{71}{\micro A}.
\end{solution}

\begin{solution}{exo:b3:electrode-kinetics:7}
(a) Reversível. (b) \omterm{def:b3:electrode-kinetics:reversibility}{Quase reversível}: a separação cresce com a velocidade de varredura.
(c) Uma etapa química segue a transferência de elétrons (EC): nas varreduras lentas, o
produto é consumido antes da varredura de volta; uma varredura rápida ultrapassa a química.
\end{solution}

\begin{solution}{exo:b3:electrode-kinetics:8}
$K_{ij}a_j/a_i = \num{2e-4} \times 0.14/\num{4.0e-3} = 0.007$: o eletrodo dá uma leitura 0.7\,\% alta.
\end{solution}

\begin{solution}{exo:b3:electrode-kinetics:9}
Mínimos quadrados: $b = \qty{0.1093}{\micro A.L.\micro g^{-1}}$, $a = \qty{0.466}{\micro A}$, $s = \qty{0.011}{\micro A}$; $c_0 = a/b =
\qty{4.26}{\micro g.L^{-1}}$, com $u(c_0) = \frac sb\sqrt{\frac1n + \frac{\bar y^2}{b^2S_{xx}}} = \qty{0.12}{\micro g.L^{-1}}$
($n = 4$, $\bar y = 1.286$, $S_{xx} = 125$).
\end{solution}

\begin{solution}{exo:b3:electrode-kinetics:10}
$Q = 0.0500 \times 1800 = \qty{90.0}{C}$; $n(\ce{Cu}) = Q/2F = \qty{4.66e-4}{mol}$, \qty{29.6}{mg}. O
resultado só depende da carga, medida a partir de uma corrente e de um tempo, e da
constante de Faraday: nenhum padrão é necessário, desde que a eletrólise seja
completa e seu rendimento faradaico seja de 100\,\%.
\end{solution}

\begin{solution}{exo:b3:electrode-kinetics:11}
$\sqrt{DFv/RT} = \sqrt{\num{1.0e-5} \times 38.92 \times 0.1} = \qty{6.24e-3}{cm.s^{-1}}$: $\Lambda = 1.6$, \omterm{def:b3:electrode-kinetics:reversibility}{quase reversível}.
A \qty{10}{V/s}, $\Lambda = 0.16$, ainda \omterm{def:b3:electrode-kinetics:reversibility}{quase reversível}, porém mais longe do limite
reversível; a onda só seria reversível abaixo de cerca de \qty{1}{mV/s}.
\end{solution}

\begin{solution}{exo:b3:electrode-kinetics:12}
$s/b = \qty{0.0838}{mM}$ e $b^2S_{xx} = \qty{203.5}{\micro A^2}$. Em $\bar y$: $0.0838\sqrt{1 + 1/7} = \qty{0.090}{mM}$.
Em \qty{18.0}{\micro A}: $0.0838\sqrt{1.143 + 83.0/203.5} = \qty{0.104}{mM}$. Três leituras em $\bar y$:
$0.0838\sqrt{1/3 + 1/7} = \qty{0.058}{mM}$.
\end{solution}

\begin{solution}{pb:b3:electrode-kinetics:1}
\textbf{1.} $\ce{C6H12O6 -> C6H10O6 + 2 H+ + 2 e-}$ e $\ce{[Fe(CN)6]^3- + e- -> [Fe(CN)6]^4-}$;
global
\[
  \ce{C6H12O6 + 2 [Fe(CN)6]^3- -> C6H10O6 + 2 [Fe(CN)6]^4- + 2 H+} .
\]
No eletrodo, $\ce{[Fe(CN)6]^4- -> [Fe(CN)6]^3- + e-}$.
\textbf{2.} Dois.
\textbf{3.} O oxigênio dissolvido no sangue é escasso e variável, e seu produto, o
peróxido de hidrogênio, só é oxidado em um potencial alto, no qual muitas outras
espécies reagem; o mediador está presente em excesso conhecido.
\textbf{4.} Para que sua concentração superficial seja nula: a corrente fica então
limitada apenas pela difusão e insensível a pequenas variações de potencial.
\textbf{5.} A corrente cai como $t^{-1/2}$; as leituras precisam ser feitas no tempo usado
na calibração.
\textbf{6.} Cottrell: $i\sqrt t = 42.0$, 41.4, 41.7, 41.6, 41.8: constante a 1\,\%.
\textbf{7.} $D = \pi(\text{inclinação}/FAc)^2 = \pi(\num{41.8e-6}/(96\,485 \times 0.0300 \times \num{1.00e-5}))^2 = \qty{6.55e-6}{cm^2.s^{-1}}$.
\textbf{8.} $\sqrt{\pi \times \num{6.55e-6} \times 5} = \qty{0.010}{cm} = \qty{100}{\micro m}$: a depleção atinge o topo da
camada por volta de \qty{5}{s}; leituras posteriores ficariam abaixo da lei de Cottrell.
\textbf{9.} A metade: \qty{9.35}{\micro A}.
\textbf{10.} A corrente de fundo: oxidação de interferentes e impurezas, carga
capacitiva, o hexacianoferrato(II) residual do próprio mediador.
\textbf{11.} Sim: os resíduos, no máximo \qty{0.095}{\micro A}, se dispersam sem tendência. Com muita glicose, o mediador, presente em
quantidade limitada, ou a \omterm{def:b3:complex-kinetics:enzyme}{enzima} se satura e a reta se curva.
\textbf{12.} $(7.10 - 0.356)/0.9189 = \qty{7.34}{mM}$.
\textbf{13.} $0.0838\sqrt{1 + 1/7 + (7.10 - 8.89)^2/203.5} = 0.0838 \times 1.076 = \qty{0.090}{mM}$.
\textbf{14.} O termo em $1/m$ (uma única leitura); leituras repetidas, ou várias tiras,
o reduzem.
\textbf{15.} $3 \times 0.077/0.919 = \qty{0.25}{mM}$.
\textbf{16.} $7.34 \times 180.16 = \qty{1322}{mg.L^{-1}} = \qty{132}{mg.dL^{-1}}$.
\textbf{17.} Ele acrescenta sua própria corrente de oxidação: o aparelho dá uma leitura alta demais.
\textbf{18.} Um \omterm{def:b3:electrode-kinetics:cv}{voltamograma} da tira sem glicose mostra uma onda anódica do ascorbato
onde o hexacianoferrato(II) é oxidado.
\textbf{19.} Menos substâncias são oxidadas em um potencial mais baixo.
\textbf{20.} A corrente do branco é subtraída da leitura do eletrodo com \omterm{def:b3:complex-kinetics:enzyme}{enzima} antes
de aplicar a reta de calibração.
\textbf{21.} $D$ cresce alguns por cento por kelvin, e a corrente como $\sqrt D$: os aparelhos
medem a temperatura e fazem a correção.
\textbf{22.} A viscosidade e a fração de hemácias do sangue alteram $D$ e a corrente;
padrões em uma matriz semelhante cancelam esses efeitos.
\textbf{23.} As variações de uma tira para outra (área do eletrodo, carga de \omterm{def:b3:complex-kinetics:enzyme}{enzima} e
de mediador), a temperatura e a composição do sangue se somam à incerteza própria da calibração.
\textbf{24.} \textbf{$c(\text{glicose}) = 7.34 \pm \qty{0.09}{mmol.L^{-1}}$} (incerteza-padrão), cerca de \qty{132}{mg.dL^{-1}}.
\end{solution}
